% Zone „Das kennst du schon“ – aus bank/kombinatorik/zone.jsonl (zusammenbau v0.1)
\einheitenkopf[zone][Kennst du schon]{Das kennst du schon}
\setcounter{aufgabe}{0}

%% TODO Titel: Ich-kann-Satz fehlt in der Bank (Platzhalter: Kettenname) – Nr. 1
%% TODO Anweisung: ein Satz über der ersten Teilaufgabe fehlt in der Bank – Nr. 1
\begin{aufgabe}{Möglichkeiten vollständig aufzählen (Liste, Tabelle, Baum ohne Wahrscheinlichkeiten), Zählprinzip „mal“ bei zwei Auswahlen, Anordnungen einer kleinen Menge auflisten und als Produkt zählen, Ergebnismenge als geordnete Paare}
\begin{teile}
\teil Tim hat drei Hosen und vier T-Shirts. Wie viele verschiedene Kombinationen aus einer Hose und einem T-Shirt kann er anziehen? \leerfeld
\teil Drei Bücher A, B und C stehen nebeneinander im Regal. Schreibe alle Reihenfolgen auf und zähle sie als Produkt. \leerfeld
\end{teile}
\end{aufgabe}

%% TODO Titel: Ich-kann-Satz fehlt in der Bank (Platzhalter: Kettenname) – Nr. 2
%% TODO Anweisung: ein Satz über der ersten Teilaufgabe fehlt in der Bank – Nr. 2
\begin{aufgabe}{Potenzen mit natürlichen Exponenten, Produkte großer Zahlen mit dem Rechner, Ergebnisse in Zehnerpotenzschreibweise lesen und angeben, die Rechnertasten für Fakultät und Binomialkoeffizient (n!, nCr; gerätabhängig)}
\begin{teile}
\teil Berechne $2^5$.
\teil Schreibe 4\,500\,000 in Zehnerpotenzschreibweise.
\end{teile}
\end{aufgabe}

%% TODO Titel: Ich-kann-Satz fehlt in der Bank (Platzhalter: Kettenname) – Nr. 3
%% TODO Anweisung: ein Satz über der ersten Teilaufgabe fehlt in der Bank – Nr. 3
\begin{aufgabe}{Brüche und Quotienten kürzen (Fakultätenquotienten, Anteile aus zwei Anzahlen), Anteil als Dezimalzahl und Prozent, Vergleich mit einer Schranke}
\begin{teile}
\teil Kürze $\frac{16}{40}$ vollständig.
\teil Kürze und berechne $\frac{8 \cdot 7 \cdot 6}{3 \cdot 2 \cdot 1}$.
\end{teile}
\end{aufgabe}

%% TODO Titel: Ich-kann-Satz fehlt in der Bank (Platzhalter: Kettenname) – Nr. 4
%% TODO Anweisung: ein Satz über der ersten Teilaufgabe fehlt in der Bank – Nr. 4
\begin{aufgabe}{Pfadregeln}
\begin{teile}
\teil Ein Würfel wird zweimal geworfen. Wie groß ist die Wahrscheinlichkeit für zwei Sechsen? $P = $ \leerfeld
\teil Ein Würfel wird zweimal geworfen. Wie groß ist die Wahrscheinlichkeit für mindestens eine Sechs? $P = $ \leerfeld
\end{teile}
\end{aufgabe}

%% TODO Titel: Ich-kann-Satz fehlt in der Bank (Platzhalter: Kettenname) – Nr. 5
%% TODO Anweisung: ein Satz über der ersten Teilaufgabe fehlt in der Bank – Nr. 5
\begin{aufgabe}{Werte einer Folge mit dem Rechner berechnen und gegen eine Schranke prüfen (beide Nachbarwerte belegen)}
\begin{teile}
\teil Ab welchem $n$ ist $2^n$ größer als 100? Nenne beide Nachbarwerte. $n = $ \leerfeld
\teil Bestimme mit dem Rechner das kleinste $n$ mit $0{,}9^n < 0{,}5$. Nenne beide Nachbarwerte. $n = $ \leerfeld
\end{teile}
\end{aufgabe}

%% TODO Titel: Ich-kann-Satz fehlt in der Bank (Platzhalter: Kettenname) – Nr. 6
%% TODO Anweisung: ein Satz über der ersten Teilaufgabe fehlt in der Bank – Nr. 6
\begin{aufgabe}{Möglichkeiten vollständig aufzählen (Liste, Tabelle, Baum ohne Wahrscheinlichkeiten), Zählprinzip „mal“ bei zwei Auswahlen, Anordnungen einer kleinen Menge auflisten und als Produkt zählen, Ergebnismenge als geordnete Paare – Fehler finden}
\begin{teile}
\teil Tom soll alle dreistelligen Zahlen aus den Ziffern 4, 5 und 6 ohne doppelte Ziffer aufschreiben. Finde den Fehler und gib die richtige Anzahl an. \rechnung{&456, \ 465, \ 546, \ 654 \\ &\text{also } 4 \text{ Zahlen}}
\end{teile}
\end{aufgabe}

%% TODO Titel: Ich-kann-Satz fehlt in der Bank (Platzhalter: Kettenname) – Nr. 7
%% TODO Anweisung: ein Satz über der ersten Teilaufgabe fehlt in der Bank – Nr. 7
\begin{aufgabe}{Möglichkeiten vollständig aufzählen (Liste, Tabelle, Baum ohne Wahrscheinlichkeiten), Zählprinzip „mal“ bei zwei Auswahlen, Anordnungen einer kleinen Menge auflisten und als Produkt zählen, Ergebnismenge als geordnete Paare – selbst rechnen}
\begin{teile}
\teil Aus den Ziffern 2, 5, 7 und 9 werden zweistellige Zahlen ohne doppelte Ziffer gebildet. Schreibe alle nach der ersten Ziffer geordnet auf. Wie viele sind es? \leerfeld
\end{teile}
\end{aufgabe}
